\documentclass[10pt]{article}

\usepackage[utf8]{inputenc}

\usepackage[T1]{fontenc}

\usepackage{amsmath}

\usepackage{amsfonts}

\usepackage{amssymb}

\usepackage[version=4]{mhchem}

\usepackage{stmaryrd}

\usepackage{bbold}

\usepackage{mathrsfs}

\usepackage{CJKutf8}

\usepackage{graphicx}

\usepackage[export]{adjustbox}

\graphicspath{ {./images/} }



\begin{document}

Множество всех С. м. $n$-го порядка представляет собой выпуклую оболочку $n^{n} \mathrm{C}$. м., составленных из нулей и единиц. Любую С. м. $P$ можно рассматривать как переходных вероятностей матрицу цепи Маркова $\xi^{P}(t)$ с дискретным временем.



Абсолютные величины собственных значений С. м. не превосходят единицы; единица является собственным значением любой С. м. Если С. м. $\boldsymbol{P}$ неразложима (цень Маркова $\xi^{P}(t)$ имеет один положительный класс состояний), то единица является простым собственным значением матрицы $P$ (т. е. имеет кратность 1 ); в общем случае кратность собственного значения 1 совпадает с qислом положительных классов цепи Маркова $\xi_{P}(t)$. Если С. м. неразложима и положительный класс состояний цепи Маркова имеет период $d$, то множество всех собственных значений матрицы $P$, как множество точек комплексной плоскости, переходит в себя при повороте на угол $2 \pi / d$. При $d=1$ С. м. $P$ и цепь Маркова $\xi^{P}(t)$ наз. н е п е р и о ди ч е с ким и.



Левые собственные векторы $\pi=\left\{\pi_{j}\right\}$ С. м. $P$ конечного порядка, соответствующие единичному собственному значению:



\[

\pi_{j}=\sum_{i} \pi_{i} p_{i j} \text { при всех } j,

\]



и удовлетворяющие условиям $\pi_{j} \geqslant 0, \sum_{j} \pi_{j}=1$, определяют стационарные распределения цепи Маркова $\xi^{p}(t)$; в случае неразложимой С. м. $P$ стационарное распределение единственно.



Если $P$ - неразложимая непериодическая С. м. конечного порядка, то существует



\[

\lim _{n \rightarrow \infty} P^{n}=\Pi

\]



где ПI - матрица, каждая строка к-рой совпадает с вектором $\pi$ (см. также Маркова цепь эргодическая; для бесконечных С. м. $P$ система уравнений (1) может не иметь ненулевых решений, удовлетворяющих условию $\sum_{j} \pi_{j}<\infty$; в этом случае матрица П-нулевая). Скорость сходимости в (2) можно оценить геометрич. прогрессией с любым показателем $\rho$, к-рый по модулю больше всех собственных значений матрицы $P$, отличных от 1.



Если $\stackrel{P}{P}=\left\|p_{i j}\right\|-$ С. м. $n$-го порядка, то любое ее собственное значение $\lambda$ удовлетворяет неравенству (cM. [3]):



\[

|\lambda-\omega| \leqslant 1-\omega, \text { где } \omega=\min _{1 \leqslant i \leqslant n} p_{i i} \text {. }

\]



Описано мноякество $M_{n}$, являющееся объединением множеств собственных значений всех С. м. $n$-го порядка (см. [4]).



C. м. $P=\left\|p_{i j}\right\|$, удовлетворяющая дополнительному условию



\[

\sum_{i} p_{i j}=1 \text { при всех } j,

\]



наз. Д в а жы с тох а с тй е с к о й м а т р иц е й. Множество дважды стохастич. матриц $n$-го порядка представляет собой выпуклую оболочку $n$ ! перестановочных матриц $n$-го порядка (т. е. дважды стохастич. матриц, составленных из нулей и единиц). Стационарное распределение конечной цепи Маркова $\xi P(t)$ с дважды стохастич. матрицей $P$ является равномерным.



Лum.: [1] г а н т м a $\mathrm{x}$ е $\mathrm{p}$ Ф. Р., Теория матриц, 3 изд., М., 1967; [2] Б е л Л м а н Р., Введение в теорию матриц, пер. с англ., М., 1969; [3] М а к у с М., М и н к Х., Обзор по тео-

рии матриц и матричных неравенств, пер. с англ., М., 1972; [4] К а р п е л е в и ч Ф. И., «Изв. АН СССР. Сер. матем.»,

1951, т. 15 , с. 361-83.

A. М. Зубков. 1951, 'T. 15 , с. 361 - 83 . выборочных функций случайного процесса. Случайный процесс $X(t)$, заданный на нек-ром множкестве $T \subseteq \mathbb{R}^{1}$, наз. с т ох а с ти че ск и не п р е р ы в ы ма этом множестве, если для любого $\varepsilon>0$ при всех $t_{0} \in T$ $\lim _{t \rightarrow t_{0}} \mathrm{P}\left\{\rho\left[X(t), \quad X\left(t_{0}\right)\right]>\varepsilon\right\}=0$



где $\rho$ - расстояние между точками в соответствующем пространстве значений процесса $X(t)$.



Лит.: [1] П р о х о р о в Ю. В., Р о з а н о в Ю. А., ТеоСТОХАСТИЧЕСКАЯ̆ М НЕРАЗЛИЧИМОСТЬ - свойство двух случайных процессов $X=\left(X_{t}(\omega)\right)_{t \geqslant 0}$ и $Y=$ $=\left(Y_{t}(\omega)\right)_{t \geqslant 0}$, означающее, что случайное множество



\[

\{X \neq Y\}=\left\{(\omega, t): X_{t}(\omega) \neq Y_{t}(\omega)\right\}

\]



является пренебрежимым, т. е. вероятность множества $\{\omega: \exists t \geqslant 0$, что $(\omega, t) \in\{X \neq Y\}\}$ равна нулю. Если $X$ и $Y$ стохастически неразличимы, то $X_{t}=Y_{t}$ (для всех $t \geqslant 0$, т. е. $X$ и $Y$ - стохастически эквивалентны). Обратное, вообще говоря, неверно, но для процессов непрерывных справа (слева) из стохастич. эквивалентности следует С. н. Лит.: [1] Д е л л а ші е р п К., Емкости и случайные про-

цессы, пер. с франц., М., 1975. СТОХАСТИЧЕСКАЯ ОГРАНИЧЕННОСТЬ, о г p aни ч ен но с т в по в е р о я т но с т и, - свойство случайного процесса $X(t), t \in T$, к-рое выражается условием: для произвольного $\varepsilon>0$ существует такое $C>0$, что при всех $t \in T$



$\mathrm{P}\{|X(t)|>C\}<\varepsilon . \quad$ A. В. Прохоров.



СТОХАСТИЧЕСКАЯ ПОСЛЕДОВАТЕЛЬНОСТЬ последовательность случайных величин $X=\left(X_{n}\right)_{n \geqslant 1}$, заданная на измеримом пространстве $(\Omega, \mathscr{F})$, с выделенным на нем неубывающим семейством $\sigma$-алгебр $\left(\mathscr{F}_{n}\right)_{n \geqslant 1}$, $\mathscr{F}_{n} \subseteq \mathscr{F}$, обладающих свойством согласованности (адаптированности): $X_{n}$ при каждом $n \geqslant 1$ \begin{CJK}{UTF8}{mj}売\end{CJK} $n$-измеримы. Для записи таких последовательностей часто используется обозначение $X=\left(X_{n}, \mathscr{F}_{n}\right)_{n \geqslant 1}$, подчерки-



\includegraphics[max width=\textwidth, center]{2023_07_09_c335dec0206a004dbf1bg-1}

примерами С. п., заданных на вероятностном пространстве $(\Omega, \mathscr{F}$, Р), являются марковские последовательности, мартингаль, семимартингальь и др. В случае непрерывного времени (когда дискретние время $n \geqslant 1$ заменяется на $t \geqslant 0$ ) соответствующая совокупность объектов $X=\left(X_{t}, \mathscr{F}_{t}\right)_{t \geqslant 0}$ наз. с т о х а с т и ч с с и м п р о п е с с о м. $\quad$ A. Н. Ширлев. СТОХАСТИЧЕСКАЯ СХОДИМОСТЬ - то же, что сходимость по вероятности.



СТОХАСТИЧЕСКАЯ ЭКВИВАЛЕНТНОСТЬ - отношение эквивалентности между случайными величинами, раздичающимися лишь на множестве нулевой вероятности. Точнее, случайные величины $X_{1}$ и $X_{2}$, заданные на одном вероятностном пространстве ( $\Omega, \mathscr{J}, \mathrm{P})$, наз. с т о х а т и ч е с ки э в в в а ле нт ным и, если $\mathrm{P}\left\{X_{1}=X_{2}\right\}=1$. В большинстве задач теории вероятностей имеют дело не с самими случайными величинами, а с классами эквивалентных случайных величин.



Случайные процессы $X_{1}(t)$ и $X_{2}(t), t \in T$, определенные на одном вероятностном пространстве, наз. с т о$\mathbf{x}$ a c т ич е с $к$ и $ә$ к и в а л е н т ны ми, если при любом $t \in T$ имеет место С. э. между соответствующими случайными величинами: $\mathrm{P}\left\{X_{1}(t)=X_{2}(t)\right\}=1$. По отношению к случайным процессам $X_{1}(t)$ и $X_{2}(t)$, у к-рых совпадают соответственные конечномерные распределения, применим также термин «С. э.» в широком смысле.



СТОХАСТИЧЕСКИИ БАЗИС А. В. Прохоров. ное пространство $(\Omega, \mathscr{F}, \mathrm{P})$ с выделенным на нем неубывающим семейством $\mathbb{F}=(\mathscr{F})_{t \geqslant 0}$ под-б-алгебр $\mathscr{F} t \subseteq \mathscr{J}$, удовлетворяющих (т. н. обычным) условиям:



\begin{enumerate}

  \item непрерывность справа, \begin{CJK}{UTF8}{mj}咭\end{CJK} $t^{=} \mathscr{T}_{t+}\left(=\underset{s>t}{=\bigcap_{s}} \mathscr{F}_{s}, t \geqslant 0\right.$,

\end{enumerate}



\end{document}